\documentclass[11pt]{article}
\usepackage[margin=1in]{geometry}
\usepackage{amsmath,amssymb,amsthm,booktabs}
\usepackage[hidelinks]{hyperref}
\newtheorem{theorem}{Theorem}
\newtheorem{lemma}[theorem]{Lemma}
\newtheorem{proposition}[theorem]{Proposition}
\theoremstyle{remark}
\newtheorem{remark}[theorem]{Remark}
\title{Stability thresholds of $N$-vortex rings on the hyperbolic plane\\
(sinh model): an expository derivation, the $\cosh^2$ form, and an arithmetic remark}
\author{Gordon Speagle\\Independent Researcher, Boise, ID, USA\\[2pt]\small Expository draft (results known: Montaldi--Tokieda 2013)}
\date{2026-10-12}
\begin{document}
\maketitle

\begin{abstract}
For $N$ identical point vortices on the hyperbolic plane, with the ``sinh'' pair interaction
$-\ln\sinh(d/2)$, a regular $N$-ring of Euclidean radius $\sqrt\xi$ in the Poincar\'e disk is spectrally
stable if and only if $(N-1)(1+\xi^2)/(1-\xi)^2 \ge \lfloor N/2\rfloor\lceil N/2\rceil/2$, and its Hessian is
positive definite on the non-rotational modes if and only if the inequality is strict.
\textbf{This criterion is known}: it is Montaldi--Tokieda~\cite{MT}, eq.~(2.11) and Theorem~2. We give a short self-contained derivation in
log-polar coordinates, in which the Hessian is diagonal in each Fourier mode. We restate the threshold as
$\cosh^2 r_*=q(N)$ with $q$ rational, and we check the thresholds against an independent 60-digit
linear-stability computation for $N=8,\dots,15$. We also record an elementary remark: $\xi_*(N)$ is an
algebraic unit exactly for $N\in\{8,9,11,15,23\}$.
\end{abstract}

\section{Model}
Poincar\'e disk $\mathbb D=\{|z|<1\}$ with metric $4|dz|^2/(1-|z|^2)^2$ (curvature $-1$) and area
(symplectic) form $\omega=\sum_k \lambda(z_k)\,dx_k\wedge dy_k$, $\lambda(z)=4/(1-|z|^2)^2$. For $N$ unit
vortices:
\begin{equation}\label{H}
H=-\sum_{j<k}\ln|z_j-z_k|+\frac{N-1}{2}\sum_k\ln(1-|z_k|^2)
=-\sum_{j<k}\ln\sinh\frac{d_{jk}}{2}+\text{const},
\end{equation}
using $\sinh(d/2)=|z-w|/\sqrt{(1-|z|^2)(1-|w|^2)}$. Rotations about $0$ are generated by
$J=\sum_k(1+|z_k|^2)/(1-|z_k|^2)$ ($=\sum_k\cosh r_k$). Hamilton's equations are $\dot x_k=\lambda^{-1}\partial_{y_k}H$ and
$\dot y_k=-\lambda^{-1}\partial_{x_k}H$.

\paragraph{Which model this is.} Up to the curvature normalisation, (\ref{H}) is the $\lambda<0$ member of
the family built on the Green's function $\log\big(|z-w|^2/(\sigma_z\sigma_w)\big)$, $\sigma=1+\lambda|z|^2$,
in Montaldi--Tokieda~\cite[(1.12)]{MT}. We call it the \emph{sinh model}.

\section{Derivation of the mode decomposition}
Write $z_k=\sqrt{\xi}\,e^{s_k+i(2\pi k/N+\theta_k)}$, so the ring is $s=\theta=0$. In these coordinates
$\omega=\sum_k\kappa(s_k)\,ds_k\wedge d\theta_k$ with $\kappa=4\xi e^{2s}/(1-\xi e^{2s})^2$, and
$\kappa_0:=\kappa(0)=4\xi/(1-\xi)^2$ at the ring.

\paragraph{Pair term.} With $\phi=2\pi(j-k)/N$, $\Delta s=s_j-s_k$, $\Delta\theta=\theta_j-\theta_k$,
\[
\ln|z_j-z_k|^2=s_j+s_k+\ln\xi+\ln\!\big(2\cosh\Delta s-2\cos(\phi+\Delta\theta)\big).
\]
Expanding to second order (sympy, \texttt{verify\_core.py}) gives no first-order $\Delta s$ term, no
$\Delta s\,\Delta\theta$ cross term, and quadratic part $(\Delta s^2-\Delta\theta^2)/D_\phi$ with
$D_\phi=2-2\cos\phi$. Hence the quadratic part of $H_{\rm pair}=-\frac12\sum_{j<k}\ln|z_j-z_k|^2$ is
$-\tfrac12 Q(s)+\tfrac12 Q(\theta)$, where $Q(x)=\sum_{j<k}(x_j-x_k)^2/D_{jk}$.

\begin{lemma}\label{fejer}
For $0\le m\le N$: $\displaystyle\sum_{l=1}^{N-1}\frac{1-\cos(2\pi ml/N)}{1-\cos(2\pi l/N)}=m(N-m)$.
\end{lemma}
\begin{proof}
Let $w=e^{2\pi il/N}$. For $l\ne0$, $\frac{1-\cos(2\pi ml/N)}{1-\cos(2\pi l/N)}=\frac{|1-w^m|^2}{|1-w|^2}
=\big|\sum_{a=0}^{m-1}w^a\big|^2=\sum_{a,b=0}^{m-1}w^{a-b}$. Summing over $l=0,\dots,N-1$ gives
$N\cdot\#\{(a,b):a\equiv b\bmod N\}=Nm$, because $|a-b|<m\le N$. Subtracting the $l=0$ term $m^2$ leaves
$m(N-m)$.
\end{proof}
(Equivalent to the identity (2.4) quoted in~\cite{MT}. Checked numerically to 50 digits for $N\le80$.)
On the Fourier mode $x_k=e^{2\pi imk/N}$, the lemma gives $Q(x)=\frac{m(N-m)}{2}\sum_k|x_k|^2$. So the
Hessian of $H_{\rm pair}$ is diagonal on mode $m$, with entry $-m(N-m)/2$ in $s$ and $+m(N-m)/2$ in $\theta$.

\paragraph{One-body terms.} $W(s)=\frac{N-1}2\ln(1-\xi e^{2s})$ and
$J_1(s)=(1+\xi e^{2s})/(1-\xi e^{2s})$ do not depend on $\theta$. The linear part $-\frac{N-1}{2}\sum s_k$ of
$H_{\rm pair}$ together with $W'(0)$ and $-\Omega J_1'(0)$ gives the equilibrium condition, hence
\[
\Omega=-\frac{(N-1)(1+\xi)(1-\xi)}{8\xi},\qquad
W''(0)-\Omega J_1''(0)=-\frac{2(N-1)\xi}{(1-\xi)^2}+\frac{(N-1)(1+\xi)^2}{(1-\xi)^2}
=\frac{(N-1)(1+\xi^2)}{(1-\xi)^2}=:C_1(\xi).
\]

\begin{proposition}[known~\cite{MT}]\label{prop:modes}
In the coordinates $(s,\theta)$, the Hessian of $F=H-\Omega J$ at the ring is block-diagonal over Fourier
modes $m=0,\dots,N-1$. On mode $m$ it is $\mathrm{diag}(a_m,b_m)$ with
\[
a_m=C_1(\xi)-\tfrac{m(N-m)}{2},\qquad b_m=\tfrac{m(N-m)}{2}.
\]
The linearised flow on mode $m$ is $\dot s=(b_m/\kappa_0)\theta$, $\dot\theta=-(a_m/\kappa_0)s$, with
eigenvalues $\sigma^2=-a_mb_m/\kappa_0^2$.
\end{proposition}
Because $b_m>0$ for $1\le m\le N-1$, and $m=0$ is the rotation direction ($b_0=0$), the ring is spectrally
stable (all eigenvalues of the linearised flow purely imaginary) iff $a_m\ge0$ for all $m$. It has a real
eigenvalue pair $\pm\sigma$, $\sigma>0$, iff some $a_m<0$. Since
$\max_m m(N-m)/2=\lfloor N/2\rfloor\lceil N/2\rceil/2=:T$, the criterion is $C_1(\xi)\ge T$. At equality the
maximising mode has $a_m=0$, $b_m>0$, so its block is nilpotent (linear growth): the ring is spectrally stable but not
linearly stable there. When
$C_1(\xi)>T$, the Hessian of $F$ is positive definite transverse to the rotation orbit, and
Montaldi--Tokieda, Theorem~2 (stated for exactly this model), gives Lyapunov stability modulo rotations. At
$C_1(\xi)=T$ the Hessian is only semidefinite; that case is not treated in this note. For $N\ge8$ it is the
bifurcation point of Montaldi--Tokieda's Conjecture~13, and the companion paper in \texttt{conj13/paper} of this
repository proves Lyapunov stability modulo rotations there. We did not
re-derive the energy--momentum step.

\section{Marginal stability, the $\cosh^2$ form}
With $\xi=\tanh^2(r/2)$ ($r$ the hyperbolic ring radius), $\cosh r=(1+\xi)/(1-\xi)$ and
$(1+\xi^2)/(1-\xi)^2=(1+\cosh^2r)/2$. So $C_1=T$ is equivalent to
\[
\cosh^2r_*=\frac{2T}{N-1}-1=q(N),\qquad
q(N)=\begin{cases}\frac{N-3}{4}&N\text{ odd},\\[2pt] \frac{(N-2)^2}{4(N-1)}&N\text{ even},\end{cases}
\qquad \xi_*=\frac{\sqrt q-1}{\sqrt q+1}.
\]
$q(N)\le1$ for $N\le7$, so every ring is spectrally stable ($N=7$: $q=1$, $r_*=0$, the planar Thomson--Havelock
degeneracy). For $N\ge8$ we have $q>1$, the ring is spectrally stable iff $r\ge r_*(N)$, with positive definite Hessian iff $r>r_*(N)$. This is
Montaldi--Tokieda's (2.11) rewritten. Their table of bifurcation values $\lambda r_0^2=-0.0627,-0.101,
-0.143,-0.172,-0.202,-0.225$ ($n=8,\dots,13$) is $-\xi_*(N)$, and their example $n=15$ gives $\xi_*=2-\sqrt3$.

\begin{center}
\begin{tabular}{rllll}\toprule
$N$ & $q(N)$ & $\xi_*$ & $\sqrt{\xi_*}$ (Euclid.) & $r_*$ (hyperbolic)\\\midrule
8 & 9/7 & $8-3\sqrt7=0.0627460668062\ldots$ & 0.250492 & 0.511875\\
9 & 3/2 & $5-2\sqrt6=0.1010205144336\ldots$ & 0.317837 & 0.658479\\
10 & 16/9 & $1/7$ & 0.377964 & 0.795365\\
11 & 2 & $3-2\sqrt2=0.1715728752538\ldots$ & 0.414214 & 0.881374\\
12 & 25/11 & $0.2024108640318\ldots$ & 0.449901 & 0.969152\\
13 & 5/2 & $0.2251482265544\ldots$ & 0.474498 & 1.031719\\
14 & 36/13 & $0.2492775954100\ldots$ & 0.499277 & 1.096685\\
15 & 3 & $2-\sqrt3=0.2679491924311\ldots$ & 0.517638 & 1.146216\\\bottomrule
\end{tabular}
\end{center}

\paragraph{Numerical check against the dynamics (NUMERICAL).}
\texttt{sinh/lin\_stab.py} codes the Cartesian vector field of (\ref{H}) directly, without the Fourier
reduction. It solves for $\Omega$ and forms the full $2N\times2N$ Jacobian by central differences at 60
digits ($h=10^{-25}$). (A) For $N=8,\dots,15$ and $\xi\in\{0.02,\,0.9\xi_*,\,1.1\xi_*,\,0.5\}$, the spectrum
matches $\{0,0\}\cup\{\pm\sqrt{-a_mb_m}/\kappa_0\}$ to within $2.5\cdot10^{-17}$, and $\Omega$ matches the
formula to $10^{-58}$. (B) Bisection on ``$\max\operatorname{Re}\sigma>10^{-15}$'' locates the threshold
within $1.2\cdot10^{-19}$ of $\xi_*(N)$ for every $N=8,\dots,15$, which is the bisection resolution.
Recorded output: \texttt{sinh/lin\_stab.out}. Coarser independent scans (sinh branch: first stable
$r_E=0.252,0.318,0.378$) agree.

\section{An arithmetic remark}
From $\xi_*=(\sqrt q-1)/(\sqrt q+1)$:
\[
\xi_*+\xi_*^{-1}=\frac{2(q+1)}{q-1}=2+\frac{4}{q-1}=:c(N),\qquad
\frac4{q-1}=\begin{cases}\dfrac{16}{N-7}&N\text{ odd},\\[6pt]
\dfrac{16(N-1)}{N^2-8N+8}&N\text{ even}.\end{cases}
\]
\begin{proposition}\label{prop:units}
For $N\ge8$ the following are equivalent: (i) $\xi_*(N)$ is an algebraic integer; (ii) $\xi_*(N)$ is an
algebraic unit; (iii) $c(N)\in\mathbb Z$; (iv) $N\in\{8,9,11,15,23\}$. In those cases
$c=16,10,6,4,3$.
\end{proposition}
\begin{proof}
$\xi_*$ is a root of $x^2-cx+1$ and $0<\xi_*<1$. (iii)$\Rightarrow$(ii): $\xi_*$ is a root of the monic
integer polynomial $x^2-cx+1$, and $\xi_*(c-\xi_*)=1$. (ii)$\Rightarrow$(i) is trivial.
(i)$\Rightarrow$(iii): if $\xi_*\in\mathbb Q$, it would be a rational algebraic integer, hence in
$\mathbb Z$, which is impossible in $(0,1)$. Otherwise its minimal polynomial over $\mathbb Q$ is
$x^2-cx+1$, which has integer coefficients when $\xi_*$ is integral, so $c\in\mathbb Z$.
(iii)$\Leftrightarrow$(iv): for odd $N$, $16/(N-7)\in\mathbb Z$ iff $(N-7)\mid16$, i.e.\ $N\in\{9,11,15,23\}$
($N=8$ is even). For even $N\ge8$, $N^2-8N+8>0$, and $N^2-8N+8-16(N-1)=N^2-24N+24>0$ for $N\ge23$, so
$0<4/(q-1)<1$ there. Checking the even $N\in[8,22]$ leaves only $N=8$, where $4/(q-1)=14$.
\end{proof}
Thus $\xi_*(8)=8-3\sqrt7=(8+3\sqrt7)^{-1}$, a root of $\xi^2-16\xi+1$. The statement ``$\xi_*$ is a unit''
does not hold in general: $\xi_*(10)=1/7$ is a root of $x^2-(50/7)x+1$, whose two roots $1/7$ and $7$ multiply to $1$, but $1/7$ is not integral.
We did not find Proposition~\ref{prop:units} in the sources we read. It is an elementary observation about
known numbers and carries no dynamical content.

\section{Remarks on other models (context only)}
\begin{itemize}
\item \emph{tanh (image) model}, pair term $-\ln\tanh(d/2)$, same area form. Coarser independent scans
give first stable $r_E\approx0.26$, $0.336$,
$0.408$ for $N=8,9,10$, for this model (see~\cite{MNG,MT}). Its thresholds are
treated in the companion note \texttt{tanh/note\_tanh.pdf}: exact algebraic numbers, certified for $N\le16$.
\item \emph{flat hard-wall disk} (Euclidean area form, images plus the Robin self-image term). Independent
scans find $N=7,8$ rings unstable at every radius
tested. Identical Green's functions do not give identical dynamics.
\end{itemize}

\section{Formalisation status}
\texttt{lean/Havelock/SinhArithmetic.lean} (Lean 4.34.1 + Mathlib) compiles with no
\texttt{sorry}. Its axioms are \texttt{propext}, \texttt{Classical.choice} and \texttt{Quot.sound}. It
proves:
$\xi+\xi^{-1}=2+4/(q-1)$ for $\xi=(\sqrt q-1)/(\sqrt q+1)$, $q>1$;
the odd and even reductions of $4/(q-1)$;
$n/d\in\mathbb Z\iff d\mid n$;
$(N-7)\mid16\iff N\in\{9,11,15,23\}$ (odd $N>7$);
$(N^2-8N+8)\mid16(N-1)\iff N=8$ (even $N\ge8$);
$c(N)\in\mathbb Z\iff N\in\{8,9,11,15,23\}$ for all $N>7$, with the values $16,10,6,4,3$;
integrality and unit property of $\xi$ when $c\in\mathbb Z$;
$\xi_*(8)=8-3\sqrt7$ and $\xi_*(8)^2-16\xi_*(8)+1=0$;
$\xi_*(10)=1/7$ and $1/7$ not integral over $\mathbb Z$.
\texttt{lean/Havelock/SinhModes.lean} (with \texttt{Sinh/Hessian.lean}, \texttt{Sinh/Threshold.lean}) derives
Proposition~\ref{prop:modes} from the Hamiltonian: the ring is critical exactly for the stated $\Omega$, and on mode $m$
the Hessian acts by $a_m$, $b_m$, $g_m=0$; and it proves the threshold criterion in both forms (spectral, with
$\ge$; definite, with $>$), equivalently $q(N)\le\cosh^2r$, with stability at every radius for $N\le7$.
Not formalised: the passage from the $2\times2$ blocks to the full $2N$-dimensional linearised flow, Lyapunov
stability, and the converse (i)$\Rightarrow$(iii) for irrational $\xi_*$.

\section{Status}
The mode decomposition and the threshold are known (Montaldi--Tokieda 2013); here they are re-derived and checked
(symbolic and 60-digit numerics). The arithmetic remark is elementary, proved, and formalised for its arithmetic part.
The text was prepared with AI assistance; no human specialist has reviewed it.

\begin{thebibliography}{9}
\bibitem{MT} J.~Montaldi, T.~Tokieda, Deformation of geometry and bifurcations of vortex rings, in
\emph{Recent Trends in Dynamical Systems}, Springer Proc.\ Math.\ Stat.\ 35 (2013) 335--370; arXiv:1210.5662.
\bibitem{MNG} J.~Montaldi, C.~Nava-Gaxiola, Point vortices on the hyperbolic plane, J.\ Math.\ Phys.\ 55
(2014) 102702; arXiv:1403.2138.
\end{thebibliography}
\end{document}
